
\documentclass{article}
\usepackage{colortbl}
\usepackage{makecell}
\usepackage{multirow}
\usepackage{supertabular}

\begin{document}

\newcounter{utterance}

\centering \large Interaction Transcript for game `adventuregame', experiment `new{-}words\_home{-}delivery\_medium\_basic\_medium', episode 7 with Qwen3{-}30B{-}A3B{-}t0.0.
\vspace{24pt}

{ \footnotesize  \setcounter{utterance}{1}
\setlength{\tabcolsep}{0pt}
\begin{supertabular}{c@{$\;$}|p{.15\linewidth}@{}p{.15\linewidth}p{.15\linewidth}p{.15\linewidth}p{.15\linewidth}p{.15\linewidth}}
    \# & \multicolumn{2}{c}{Player} && \multicolumn{2}{c}{Game Master} \\
    \hline

    \theutterance \stepcounter{utterance}  
    & & & \multicolumn{4}{p{0.6\linewidth}}{
        \cellcolor[rgb]{0.9,0.9,0.9}{
            \makecell[{{p{\linewidth}}}]{
                \texttt{\tiny{[P1$\langle$GM]}}
                \texttt{You are playing a text adventure game. I will describe what you can perceive in the game. You write the single action you want to take in the game starting with >. Only reply with an action.} \\
\texttt{For example:} \\
\texttt{> open cupboard} \\
\\ 
\texttt{In addition to common actions, you can strood, eucty and pyell.} \\
\\ 
\texttt{Your goal for this game is: Put the oroust in the inchan, the book on the aciong and the plate on the aciong.} \\
\\ 
\texttt{Once you have achieved your goal, write "> done" to end the game.} \\
\\ 
\texttt{You are in a encal now. There are a book, a oroust, a bed and a wardrobe. The book is on the floor. The oroust is on the bed. The wardrobe is closed. There is a passage to a hallway here.} \\
            }
        }
    }
     \\ \\

    \theutterance \stepcounter{utterance}  
    & \multicolumn{4}{p{0.6\linewidth}}{
        \cellcolor[rgb]{0.561,0.737,0.561}{
            \makecell[{{p{\linewidth}}}]{
                \texttt{\tiny{[P1$\rangle$GM]}}
                \texttt{<think>} \\
\texttt{Okay, let's see. I need to put the oroust in the inchan, the book on the aciong, and the plate on the aciong. But first, I have to figure out where things are. The user is in an encal with a book, oroust, bed, wardrobe. The book is on the floor, oroust is on the bed. The wardrobe is closed. There's a passage to a hallway.} \\
\\ 
\texttt{Wait, the goal mentions "inchan" and "aciong". Maybe those are specific places. But the current description doesn't mention them. So maybe I need to go to the hallway to find them? Or perhaps the wardrobe is the inchan? Not sure. Let me check the items again.} \\
\\ 
\texttt{The oroust is on the bed. So I need to pick it up. Then put it in the inchan. But where is the inchan? Maybe the wardrobe? If the wardrobe is closed, maybe I need to open it first. The user can open the wardrobe. Then put the oroust inside. But the goal says "in the inchan", not sure. Alternatively, maybe the inchan is a different object. But the current description doesn't mention it. So maybe I need to go to the hallway to find the inchan. But the user hasn't mentioned that yet. Let me think.} \\
\\ 
\texttt{The user's first action should be to take the oroust from the bed. Then maybe go to the} \\
            }
        }
    }
    & & \\ \\

    \theutterance \stepcounter{utterance}  
    & & & \multicolumn{2}{p{0.3\linewidth}}{
        \cellcolor[rgb]{0.9,0.9,0.9}{
            \makecell[{{p{\linewidth}}}]{
                \texttt{\tiny{[GM$|$GM]}}
                \texttt{Round: 0} \\
            }
        }
    }
    & & \\ \\

    \theutterance \stepcounter{utterance}  
    & & & \multicolumn{2}{p{0.3\linewidth}}{
        \cellcolor[rgb]{0.9,0.9,0.9}{
            \makecell[{{p{\linewidth}}}]{
                \texttt{\tiny{[GM$|$GM]}}
                \texttt{command\_tag\_missing} \\
            }
        }
    }
    & & \\ \\

    \theutterance \stepcounter{utterance}  
    & & & \multicolumn{2}{p{0.3\linewidth}}{
        \cellcolor[rgb]{0.9,0.9,0.9}{
            \makecell[{{p{\linewidth}}}]{
                \texttt{\tiny{[GM$|$GM]}}
                {'goal_states_achieved': [], 'game_successfully_finished': False}
            }
        }
    }
    & & \\ \\

\end{supertabular}
}

\end{document}
